\section{Method}

This section describes the final pipeline: baselines, long-context encoding, contrastive supervised fine-tuning, dynamic negative sampling, domain-adaptive pretraining, scope filtering, learning to rank, and post-processing. We use the abbreviations DAPT, SFT, and LTR throughout.

Figure~\ref{fig:task1-system-pipeline} summarizes the final Task~1 pipeline discussed in this section, separating train/valid preprocessing, model training, and test-time prediction.

\begin{figure*}[t]
  \centering
  \includegraphics[width=\textwidth]{../shared/figures/task1_system_pipeline_en.pdf}
  \caption{Overview of the final Task~1 pipeline used in this paper. The diagram highlights the train/valid preprocessing assets, the training loop with dynamic negative sampling, and the test-time prediction path that combines lexical scoring, dense scoring, LightGBM reranking, and final post-processing.}
  \Description{A flowchart of the final Task 1 system pipeline with three sections: train-valid setup, training only, and test-time prediction. It shows raw training cases being cleaned, split, and indexed; contrastive training data and dynamic negative sampling feeding contrastive supervised fine-tuning and LightGBM reranker training; and test assets flowing through lexical scoring, dense scoring, learning-to-rank reranking, and final post-processing.}
  \label{fig:task1-system-pipeline}
\end{figure*}

\subsection{BM25 and SAILER Baselines}

We keep BM25 and SAILER as the main baselines. BM25 represents a strong lexical retrieval baseline, while SAILER corresponds to the legal encoder architecture proposed by THUIR@COLIEE 2023 \cite{li2023thuir}. In our experiments, SAILER is not retrained. We directly use the released model for embedding and similarity scoring.\footnote{\url{https://huggingface.co/CSHaitao/SAILER_en_finetune}} This setting is relatively favorable to SAILER because the released fine-tuned model may already have seen data overlapping with, or at least very close to, our validation distribution.

Our lexical features mainly come from BM25 and query likelihood with Dirichlet smoothing (QLD). For a query $q$ and document $d$, BM25 is
\begin{equation}
\mathrm{BM25}(q,d)=\sum_{t \in q}\mathrm{IDF}(t)\cdot
\frac{f(t,d)(k_1+1)}{f(t,d)+k_1\left(1-b+b\frac{|d|}{\mathrm{avgdl}}\right)},
\end{equation}
where $f(t,d)$ is the term frequency of $t$ in $d$, $|d|$ is document length, and $\mathrm{avgdl}$ is the average document length in the corpus. QLD is
\begin{equation}
\mathrm{QLD}(q,d)=\sum_{t \in q} c(t,q)\log \frac{c(t,d)+\mu p(t \mid C)}{|d|+\mu},
\end{equation}
where $c(t,q)$ and $c(t,d)$ are term counts and $p(t \mid C)$ is the collection language model.

\subsection{Long-Context Encoding with ModernBERT}

\begin{sloppypar}
To test whether truncation is the key bottleneck, we use Answer.AI's ModernBERT-base\footnote{\url{https://huggingface.co/answerdotai/ModernBERT-base}} encoder. ModernBERT natively supports 8192 tokens and is therefore better suited than SAILER for long judgments in principle \cite{warner2024modernbert}. In practice, however, if one contrastive batch contains one query, one positive, and 15 negatives, then 17 long documents must be encoded in one update step. On a single RTX 4080 SUPER 16GB, this makes 8192-token training infeasible, so all formal experiments use 4096 tokens instead.
\end{sloppypar}

For each judgment, we use the encoder's leading `[CLS]` vector as the dense retrieval representation, followed in the implementation by a projector and L2 normalization. All dense rankings in this paper are then computed with cosine similarity.

\subsection{Contrastive SFT and Static Hard Negatives}

We formulate legal case retrieval as a contrastive query--positive--negative learning problem. Each contrastive sample contains one query case, one noticed case, and 15 negative judgments. Let the $i$-th training triple contain query $q_i$, positive $p_i$, negative set $\mathcal{N}_i$, similarity function $s(\cdot,\cdot)$, and temperature $\tau$. The contrastive SFT loss is
\begin{equation}
\mathcal{L}_i=
-\log
\frac{\exp\left(s(q_i,p_i)/\tau\right)}
{\exp\left(s(q_i,p_i)/\tau\right)+\sum_{n \in \mathcal{N}_i}\exp\left(s(q_i,n)/\tau\right)}.
\end{equation}
Although each similarity score is computed pairwise, the loss jointly normalizes one positive against 15 negatives in a single softmax. In that sense, our dual-encoder contrastive training is closer to listwise learning than to classical pairwise ranking.

In early experiments, following a THUIR-style setup, we first build a hard-negative pool from BM25 top-100 candidates and then randomly sample 15 negatives per update. This is simple to implement, but the practical difficulty of the negatives is fixed entirely by lexical ranking.

\subsection{Dynamic Negative Sampling}

Our most important methodological change is replacing static hard negatives with dynamic negative sampling. At the beginning of each epoch, we compute the current model's similarity scores for query--candidate pairs and turn these scores into sampling probabilities for negatives. If $\mathcal{C}_i$ is the negative pool for query $q_i$ at epoch $e$, then the sampling probability is
\begin{equation}
\pi_e(n \mid q_i)=
\frac{\exp\left(\alpha s_e(q_i,n)\right)}
{\sum_{m \in \mathcal{C}_i}\exp\left(\alpha s_e(q_i,m)\right)},
\quad n \in \mathcal{C}_i,
\end{equation}
where $s_e$ is the current epoch similarity score and $\alpha$ controls how peaked the distribution is. In our implementation, the sampling temperature is set to 1.0, which corresponds to $\alpha=1$.
This keeps negative difficulty synchronized with the current model state rather than freezing it at a BM25-defined lexical notion of hardness.

\subsection{Domain-Adaptive Pretraining}

Beyond task-specific SFT, we apply DAPT on the Common Pile CaseLaw Access Project subset so that the encoder better adapts to legal sentences, citation patterns, and procedural language. Because of the current GPU budget, this stage runs for about 60 hours and still covers only 0.384 epoch, so we treat it as an under-saturated but useful adaptation step.

\subsection{Year-Based Scope Filtering}

We apply year-based scope filtering both during training data construction and during inference. The legal intuition is simple: a judge may cite prior cases, but not future ones. For each query, we therefore extract a decision year and restrict the negative pool and inference candidates to legally plausible older cases. Importantly, the current implementation does not discard gold positives when they fall outside the heuristic scope; such positives are still kept as supervised positive pairs.

During training, scope filtering prevents future judgments from being mixed into the negative set. During inference, we build a query-specific \texttt{query\_candidate\_scope} index offline and reuse it during similarity computation and reranking, so each query scores only candidates inside its own scope rather than the full corpus. In the shared scope index used here, the implementation rule is \texttt{candidate\_year <= query\_year + 1}, with a conservative one-year slack to absorb year-extraction noise.

\subsection{Learning to Rank (LTR)}

Following THUIR's learning-to-rank strategy in WSDM Cup 2023, web search ranking, and NTCIR-16 WWW-4 \cite{chen2023wsdmcup,li2023websearch,yang2023ntcir}, we integrate dense and lexical signals into a final ranking model. For each query--candidate pair we build 28 numeric features that can be grouped into seven families: lexical scores, dense scores, ranking positions, length statistics, placeholder statistics, year difference, and similarities computed from aggregated embeddings over 1--3 document chunks. \texttt{query\_id} and \texttt{candidate\_id} are retained only as identifiers and are not used as training features.

The ranking model is LightGBM's \texttt{LGBMRanker}, trained with the \texttt{lambdarank} objective and \texttt{ndcg} evaluation metric. Early stopping is performed on the validation set, and the best validation model is used at test time.

\subsection{Post-processing}

After query-wise reranking, we apply two kinds of post-processing. The first is common legal filtering, including query-specific year-scope filtering and self-document removal; the latter removes only the exact same case where \texttt{candidate\_id=query\_id}. The second is cut-off search on the validation set. Instead of always keeping the top 5 ranked items, we compare three strategies. \texttt{fixed\_topk} keeps the top $k$ items. \texttt{largest\_gap\_cutoff} finds the largest score drop within the top $N$ candidates and uses it as the truncation point. \texttt{ratio\_cutoff} keeps candidates whose score remains above a fixed ratio of the top score. The best parameters and results are reported in Section~7.
